\documentclass[12pt]{article}
\usepackage[utf8]{inputenc}
\usepackage[T1]{fontenc}
\usepackage{hyperref}
\title{Waanyi Language Resources for K-12 Curriculum Development}
\author{Compiled by Claude Sonnet 5.5 (claude-sonnet-5-5) for Andr\'as Kornai}
\date{}
\begin{document}
\maketitle

\section{Introduction}
Waanyi (also Wanyi, Wanji, Waanji; ISO 639-3 wny) is a Garrwan language of the lower Gulf area of northern Queensland, Australia. Wikipedia reports 16 speakers in the 2016 Australian census, down from 40 in 2011, so it may have a few speakers. Little material turned up, and I found no Waanyi dictionary.

\section{Primary Educational Resources}
No primers, textbooks or school programmes for Waanyi were found in the search results.

\section{Linguistic and Reference Materials}
\begin{itemize}
\item Wikipedia, Waanyi language: \url{https://en.wikipedia.org/wiki/Waanyi_language}
\item AIATSIS AUSTLANG entry for Waanyi: \url{https://collection.aiatsis.gov.au/austlang/language/c23}
\end{itemize}
Gavan Breen's chapter ``Wanyi and Garrwa comparative data'' (Pacific Linguistics 552, 2003) and Ilana Mushin's \emph{A Grammar of (Western) Garrwa} (Pacific Linguistics 637, 2012) were mentioned in search summaries. The Mushin grammar covers the related Garrwa language, not Waanyi. I saw no URL I could match to either work, so neither is in the json list.

\section{Religious and Community Resources}
No religious or community organisation resources for Waanyi were found in the search results.

\section{Digital and Audio Resources}
\begin{itemize}
\item Endangered Languages Project, Waanyi: \url{https://endangeredlanguages.com/lang/4190}
\end{itemize}
I found no audio or video resources.

\section{Curriculum Development Context}
Waanyi is not known to be taught in schools. A revitalization curriculum would depend on the Waanyi community and on collecting and checking archival material through AIATSIS and similar holdings. The published sources found are limited.
\end{document}